% Ética, sesgos y limitaciones
\section{Ética e implicaciones}

\begin{frame}
\frametitle{Sesgos del dato}
\begin{itemize}
    \item \textbf{Autoreporte y deseabilidad social}: el modelo aprende propensión a \emph{reportar}, no a \emph{consumir}. Subreporte correlacionado con percepción de riesgo legal.
    \item \textbf{Sesgo de selección por precio}: la submuestra con precio ($n\approx 755$) está sesgada hacia personas más expuestas al mercado.
    \item \textbf{Sesgo de medición en Y}: discrepancia documentada \texttt{consumo\_12m} vs \texttt{K\_03}. Mitigada con análisis de robustez.
\end{itemize}
\end{frame}

\begin{frame}
\frametitle{Riesgos del uso del modelo}
\begin{enumerate}
    \item \textbf{Perfilado individual}: AUC $\approx 0.73$ implica solapamiento alto entre grupos. \textit{No usar} para decisiones sobre personas (policial, laboral, escolar).
    \item \textbf{Asociaciones $\neq$ causalidad}: edad y sexo correlacionan con factores no observados (red social, entorno).
    \item \textbf{Estigma reforzado}: mapas y perfiles públicos refuerzan estereotipos. Privilegiar lecturas agregadas.
\end{enumerate}
\end{frame}

\begin{frame}
\frametitle{Uso responsable propuesto}
\textbf{Sí} (uso agregado):
\begin{itemize}
    \item Estimar tamaño de mercado bajo escenarios regulatorios.
    \item Diseño tributario simulando bases de contribuyentes por estratos amplios.
    \item Priorizar segmentos para política preventiva.
\end{itemize}
\vspace{0.3cm}
\textbf{No} (uso individual):
\begin{itemize}
    \item Decisiones sobre personas específicas.
    \item Focalización geográfica fina sin análisis de equidad.
    \item Predicciones puntuales sin intervalos de incertidumbre.
\end{itemize}
\end{frame}

\begin{frame}
\frametitle{Implicaciones de política}
\textbf{Diseño tributario}:
\begin{itemize}
    \item Base concentrada en cohortes jóvenes-masculinas.
    \item Elasticidad-precio probablemente alta $\Rightarrow$ impuesto excesivo desplaza al mercado ilegal.
    \item Recaudo creíble requiere acoplar este modelo con una elasticidad estimada limpia.
\end{itemize}
\vspace{0.3cm}
\textbf{Política preventiva}:
\begin{itemize}
    \item Focalizar en 18--25 años, diferenciada por sexo.
    \item Aprovechar la sensibilidad de la cohorte a intervenciones tempranas.
\end{itemize}
\end{frame}
